\documentclass{article}
%  ╔══════════════════════════════════════════════════════════╗
%  ║                         Title                            ║
%  ╚══════════════════════════════════════════════════════════╝
%NOTE: I want \title,\author to be near top, and this \renewcommand{\maketitle} fails if these are above it.
\usepackage[utf8]{inputenc}         % Used to compile any UTF-8 Character (and supports Unicode)
\usepackage{titling}

\renewcommand{\maketitle}{
    \quad
    \vspace{1em}
  \begin{center}
      \LARGE\bfseries \thetitle \par
      \vspace{0.6em}
      \large
  \end{center}
    \vspace{1em}
}

\title{Fourier Transforms}
\author{Nathanael Chwojko-Srawley}
\date{\today}

%  ╔══════════════════════════════════════════════════════════╗
%  ║                         packages                         ║
%  ╚══════════════════════════════════════════════════════════╝

\usepackage{inputenc}         % Used to compile any UTF-8 Character (and supports Unicode)
\usepackage{extarrows}         % Used to compile any UTF-8 Character (and supports Unicode)
\usepackage{makeidx}
\usepackage{stackengine}
\usepackage{scalerel}
\usepackage{amsmath, amssymb, bm} % Used to write better mathematical expressions (ex. \mathbb{R})
\usepackage{stackrel}         % for left-right arrows, staking symbol above and bellow
\usepackage{caption}
\usepackage{subcaption}
\usepackage{mathrsfs}
\usepackage{euscript}
\usepackage[font=itshape]{quoting}
\usepackage{mathtools}
\usepackage{enumitem}
\usepackage{amsthm}           % Used to create theorem enviroments.
\usepackage{graphicx}         % Let's you use images in LaTeX; ex the \includegraphics command.
\usepackage{hhline}
\usepackage{wasysym}          % emoji's!
\usepackage{dsfont}
\usepackage{float}            % package with ``H'' for figures and tables, ex. Images, tables, etc
\usepackage{setspace}         % More flexibility on spacing in document.
\usepackage{commutative-diagrams}
\usepackage{tikz}             % for commutative diagrams https://tex.stackexchange.com/questions/419221/how-to-do-the-pushout-with-universal-property
\usepackage{tikz-cd}          % for commutative diagrams https://tex.stackexchange.com/questions/419221/how-to-do-the-pushout-with-universal-property
\usepackage{longtable}        % create tables that extend past one page
\usepackage{microtype}        % makes nice text. Fixes some under-/over- filled box problems.
\usepackage{xcolor}           % Colour text.
\usepackage{wrapfig}          % To wrap text around figure
\usepackage{listings}         % Create nice colour code blocks (customized in preamble)
\usepackage{array}            % \begin{table}{>{\em}c c} -> 1st column italic (bfseries for bold)
\usepackage{blindtext}
\usepackage[a4paper, total={6in, 8in}]{geometry}
\usepackage{fancyhdr}
\usepackage{titlesec}
\usepackage[answerdelayed]{exercise}
\usepackage{etoolbox}          % for Dynkin diagrams
\usepackage{dynkin-diagrams}   % for Dynkin diagrams
\usepackage{pgfplots}
\usepackage[most]{tcolorbox} 	      % Makes nice boxes for thms, cor, prop, or anything else.
\usepackage{thmtools}         % Customize theorem environment (in preamble).
\usepackage{hyperref}
\usepackage{cleveref}

% ╔═════════════════════════════════════════════════════╗
% ║           for figures via inkspace                  ║
% ╚═════════════════════════════════════════════════════╝

% to import figures via: https://github.com/gillescastel/inkscape-figures
\usepackage{import}
\usepackage{pdfpages}
\usepackage{transparent}


% This command is the ONLY command imported via the packages snippet, think of it as a tiny package written out
\newcommand{\incfig}[2][1]{%
    \def\svgwidth{#1\columnwidth}
    \import{./figures/}{#2.pdf_tex}
}

% ╔═══════════════════════════════════════════════════════════╗
% ║                 bibliography and citing                   ║
% ╚═══════════════════════════════════════════════════════════╝

\usepackage{csquotes}
% \usepackage[backend=biber,citestyle=alphabetic,bibstyle=authortitle]{biblatex}
\usepackage{biblatex}

\addbibresource{bibliography.bib}

% ╔═══════════════════════════════════════════════════════════╗
% ║                    colours (colors)                       ║
% ╚═══════════════════════════════════════════════════════════╝

\definecolor{dkgreen}{rgb}{0,0.6,0}
\definecolor{gray}{rgb}{0.5,0.5,0.5}
\definecolor{mauve}{rgb}{0.58,0,0.82}
\definecolor{lightyellow}{rgb}{1,1,0.6}
\definecolor{reddish}{HTML}{A01A3A}

%  ╔══════════════════════════════════════════════════════════╗
%  ║                       book format                        ║
%  ╚══════════════════════════════════════════════════════════╝

\pagestyle{fancy}

\setlength\headheight{20pt}

\renewcommand{\sectionmark}[1]{\markright{\thesection.\ #1}}

\fancyfoot[C,CO]{\textcolor{gray}{Nathanael Chwojko-Srawley}}
\fancyfoot[R,RO]{\thepage}{}

% Create a new page style for the first page
\fancypagestyle{firstpage}{
  \fancyhf{} % Clear header and footer
  \fancyhead[L]{\theauthor}
  \fancyhead[R]{\today} % Date in the top-right
  \fancyfoot[C]{\textcolor{gray}{Nathanael Chwojko-Srawley}}
  \fancyfoot[R]{\thepage}
}

% Apply the first page style conditionally
\AtBeginDocument{\thispagestyle{firstpage}}

%have format the title command
\newcommand{\chapter}[1]{%
  \clearpage
  \vspace*{30pt} % Space before the chapter title
  {\normalfont\huge\bfseries #1\par} % Chapter title formatting
  \vspace{20pt} % Space after the chapter title
}


\setlength{\parindent}{0pt} % don't like indentation infront of paragraphs
\setlength{\parskip}{1ex}   %so that there is still space between paragarphs


%  ╔══════════════════════════════════════════════════════════╗
%  ║                       Shortcuts                          ║
%  ╚══════════════════════════════════════════════════════════╝

%\ExerciseLevelInToc

% I like original mathcal better, but need lowercase.
\DeclareFontFamily{OT1}{pzc}{}
\DeclareFontShape{OT1}{pzc}{m}{it}{<-> s * [1.10] pzcmi7t}{}
\DeclareMathAlphabet{\mathpzc}{OT1}{pzc}{m}{it}


% Custom Commands
\newcommand{\A}{{\mathbb{A}}}
\newcommand{\C}{{\mathbb{C}}}
\newcommand{\D}{\mathbb{D}}
\newcommand{\F}{{\mathbb{F}}}
\newcommand{\N}{{\mathbb{N}}}
\newcommand{\Q}{{\mathbb{Q}}}
\newcommand{\R}{{\mathbb{R}}}
\newcommand{\T}{{\mathbb{T}}}
\newcommand{\V}{{\mathbb{V}}}
\newcommand{\Z}{{\mathbb{Z}}}
\renewcommand{\H}{{\mathbb{H}}}
\renewcommand{\P}{{\mathbb{P}}}


\newcommand{\SA}{\mathscr{A}}
\newcommand{\SH}{\mathscr{H}}
\newcommand{\SF}{\mathscr{F}}
\newcommand{\SG}{\mathscr{G}}
\newcommand{\SK}{\mathscr{K}}
\newcommand{\SN}{\mathscr{N}}
\newcommand{\SQ}{\mathscr{Q}}
\newcommand{\SR}{\mathscr{R}}
\newcommand{\SV}{\mathscr{V}}
\newcommand{\SZ}{\mathscr{Z}}
\newcommand{\SC}{\mathscr{C}}
\newcommand{\SP}{\mathscr{P}}
\newcommand{\SI}{\mathscr{I}}


\newcommand{\TA}{\mathtt{A}}
\newcommand{\TF}{\mathtt{F}}
\newcommand{\TG}{\mathtt{G}}
\newcommand{\TN}{\mathtt{N}}
\newcommand{\TQ}{\mathtt{Q}}
\newcommand{\TR}{\mathtt{R}}
\newcommand{\TV}{\mathtt{V}}
\newcommand{\TZ}{\mathtt{Z}}
\newcommand{\TC}{\mathtt{C}}

\newcommand{\CA}{{\mathcal{A}}}
\newcommand{\CB}{{\mathcal{B}}}
\newcommand{\CE}{{\mathcal{E}}}
\newcommand{\CC}{{\mathcal{C}}}
\newcommand{\CD}{{\mathcal{D}}}
\newcommand{\CF}{{\mathcal{F}}}
\newcommand{\CG}{{\mathcal{G}}}
\newcommand{\CH}{{\mathcal{H}}}
\newcommand{\CI}{{\mathcal{I}}}
\newcommand{\CL}{{\mathcal{L}}}
\newcommand{\CM}{{\mathcal{M}}}
\newcommand{\CN}{{\mathcal{N}}}
\newcommand{\CO}{{\mathcal{O}}}
\newcommand{\co}{{\mathpzc{o}}}
\newcommand{\CQ}{{\mathcal{Q}}}
\newcommand{\CR}{{\mathcal{R}}}
\newcommand{\CU}{{\mathcal{U}}}
\newcommand{\CK}{{\mathcal{K}}}
\newcommand{\CS}{{\mathcal{S}}}
\newcommand{\CT}{{\mathcal{T}}}
\newcommand{\CV}{{\mathcal{V}}}
\newcommand{\CZ}{{\mathcal{Z}}}
% EXCPETION:
\newcommand{\CP}{\mathbb{C}\mathrm{P}}
\newcommand{\RP}{\mathbb{R}\mathrm{P}}

\newcommand{\ff}{{\mathscr{F}}}

\newcommand{\fa}{{\mathfrak{a}}}
\newcommand{\fp}{{\mathfrak{p}}}
\newcommand{\fg}{{\mathfrak{g}}}
\newcommand{\fb}{{\mathfrak{b}}}
\newcommand{\fc}{{\mathfrak{c}}}
\newcommand{\fd}{{\mathfrak{d}}}
\newcommand{\fm}{{\mathfrak{m}}}
\newcommand{\fn}{{\mathfrak{n}}}
\newcommand{\fo}{{\mathfrak{o}}}
\newcommand{\fq}{{\mathfrak{q}}}
\newcommand{\FA}{{\mathfrak{A}}}
\newcommand{\FB}{{\mathfrak{B}}}
\newcommand{\FP}{{\mathfrak{P}}}
\newcommand{\FK}{{\mathfrak{K}}}
\newcommand{\FC}{{\mathfrak{C}}}
\newcommand{\FD}{{\mathfrak{D}}}
\newcommand{\FM}{{\mathfrak{M}}}
\newcommand{\FN}{{\mathfrak{N}}}
\newcommand{\FO}{{\mathfrak{O}}}


\newcommand{\EB}{\EuScript{B}}
\newcommand{\EC}{\EuScript{C}}
\newcommand{\EF}{\EuScript{F}}
\newcommand{\EH}{\EuScript{H}}
\newcommand{\EL}{\EuScript{L}}
\newcommand{\EM}{\EuScript{M}}
\newcommand{\EN}{\EuScript{N}}
\newcommand{\EO}{\EuScript{O}}

\newcommand{\sse}{\subseteq}
\newcommand{\nsse}{\not\subseteq}
\newcommand{\ssne}{\subsetneq}
\newcommand{\qn}{\quad\newline}
\newcommand{\wto}{\rightharpoonup}
\newcommand{\placeholder}{\_\_\_\_\_}
\newcommand{\subg}{\leqslant}
%\newcommand{\subgne}{\lneqslant}
\newcommand{\supg}{\geqslant}
\newcommand{\nsubg}{\trianglelefteq}
\newcommand{\csubg}{\blacktriangleleft}
\newcommand{\nsupg}{\trianglerighteq}
\newcommand{\uniflim}{\operatorname{unif lim}}
\newcommand{\isosubg}{\lesssim}
\renewcommand{\mid}{\big|}
\renewcommand{\phi}{\varphi}

\newcommand{\oln}[1]{{\overline{#1}}}
\renewcommand{\bf}[1]{\textbf{#1}}

\renewcommand{\mod}{\,\operatorname{mod}\,}

\newcommand{\ev}{{\operatorname{ev}}}
\newcommand{\ab}{{\operatorname{ab}}}
\newcommand{\an}{{\operatorname{an}}}
\newcommand{\tr}{{\operatorname{tr}}}
\newcommand{\nm}{{\operatorname{Nm}}}
\newcommand{\op}{{\operatorname{op}}}
\newcommand{\ad}{{\operatorname{ad}}}
\newcommand{\rk}{{\operatorname{rk}}}
\newcommand{\im}{{\operatorname{im}}}
\newcommand{\id}{{\operatorname{id}}}
\newcommand{\ch}{{\operatorname{ch}}}
\newcommand{\sh}{{\operatorname{sh}}}
\newcommand{\Cl}{{\operatorname{Cl}\,}}
\newcommand{\CaCl}{{\operatorname{CaCl}\,}}
\newcommand{\SL}{{\operatorname{SL}}}
\newcommand{\GL}{{\operatorname{GL}}}
\newcommand{\aff}{{\textbf{aff}}}
\newcommand{\Sh}{{\textbf{Sh}}}
\newcommand{\Ab}{{\textbf{Ab}}}
\newcommand{\Set}{{\textbf{Set}}}
\newcommand{\Mod}{{\textbf{Mod}}}
\newcommand{\QCoh}{{\textbf{QCoh}}}
\newcommand{\Coh}{{\textbf{Coh}}}
\newcommand{\PSh}{{\textbf{PSh}}}
\newcommand{\Sch}{{\textbf{Sch}}}
\newcommand{\Top}{{\textbf{Top}}}
\newcommand{\RgSp}{{\textbf{RgSp}}}
\newcommand{\Grp}{{\textbf{Grp}}}
\newcommand{\Ring}{{\textbf{Ring}}}
\newcommand{\CRing}{{\textbf{CRing}}}
\newcommand{\Alg}{{\textbf{Alg}}}
\newcommand{\dAlg}{{\textbf{-Alg}}}
\newcommand{\rad}{{\operatorname{rad}}}
\newcommand{\Der}{{\operatorname{Der}}}
\newcommand{\res}{{\operatorname{res}}}
\newcommand{\pre}{{\operatorname{pre}}}
\newcommand{\Gal}{{\operatorname{Gal}}}
\newcommand{\Rep}{{\operatorname{Rep}}}
\newcommand{\Aut}{{\operatorname{Aut}}}
\newcommand{\Sym}{{\operatorname{Sym}}}
\newcommand{\Hom}{{\operatorname{Hom}}}
\newcommand{\Emb}{{\operatorname{Emb}}}
\newcommand{\Ext}{{\operatorname{Ext}}}
\newcommand{\SExt}{{\mathcal{E}\emph{xt}}}
\newcommand{\Pic}{{\operatorname{Pic}\,}}
\newcommand{\End}{{\operatorname{End}}}
\newcommand{\Inn}{{\operatorname{Inn}}}
\newcommand{\Ind}{{\operatorname{Ind}}}
\newcommand{\Out}{{\operatorname{Out}}}
\newcommand{\Jac}{{\operatorname{Jac}}}
\newcommand{\Nil}{{\operatorname{Nil}}}
\newcommand{\CDiv}{{\operatorname{CDiv}\,}}
\newcommand{\WDiv}{{\operatorname{WDiv}\,}}
\renewcommand{\div}{{\operatorname{div}\,}}
\newcommand{\ord}{{\operatorname{ord}\,}}
\newcommand{\red}{{\operatorname{red}}}
\newcommand{\Tor}{{\operatorname{Tor}}}
\newcommand{\Rat}{{\operatorname{Rat}}}
\newcommand{\Ann}{{\operatorname{Ann}}}
\newcommand{\Syl}{{\operatorname{Syl}}}
\newcommand{\SHom}{{\emph{Hom}}}
\newcommand{\supp}{{\operatorname{supp}}}
\newcommand{\sgn}{{\operatorname{sgn}}}
\newcommand{\conj}{{\operatorname{conj}}}
\newcommand{\kdim}{{\operatorname{kdim}}}
\newcommand{\coker}{{\operatorname{coker}}}
\newcommand{\Obj}{{\operatorname{Obj}}}
\newcommand{\Mor}{{\operatorname{Mor}}}
\newcommand{\Frac}{{\operatorname{Frac}}}
\newcommand{\Spec}{{\operatorname{Spec}\,}}
\newcommand{\Proj}{{\operatorname{Proj}\,}}
\newcommand{\mSpec}{{\operatorname{mSpec}\,}}
\newcommand{\smooth}{{\operatorname{smooth}}}
\newcommand{\Tens}[1]{#1\otimes \_\_}
\newcommand{\Char}{{\operatorname{char}}}
\newcommand{\Span}{{\operatorname{span}}}
\newcommand{\codim}{{\operatorname{codim}}}
\newcommand{\esssup}{{\operatorname{ess sup}}}

\newcommand{\Lim}{\varprojlim}
\newcommand{\colim}{\varinjlim}
\newcommand{\coLim}{\varinjlim}
\newcommand{\Dlim}{\lim\limits_{\longrightarrow}}
\newcommand{\coDlim}{\lim\limits_{\longleftarrow}}
\newcommand{\fHom}[1]{\operatorname{Hom}(#1, \_\_)}
\newcommand{\cofHom}[1]{\operatorname{Hom}(\_\_, #1)}
\newcommand{\trdeg}{\operatorname{trdeg}}
\newcommand{\acton}{\curvearrowright}
\newcommand{\actson}{\curvearrowright}

\newcommand\cH{\check{\mathrm{H}}}

\renewcommand{\Re}{{\operatorname{Re}}}
\renewcommand{\Im}{{\operatorname{Im}}}

%for saturated ideals in the localization section
\newcommand{\LIm}{{}^{\iota} }
\newcommand{\LPre}{{}^{\pi} }

% Arrows
\newcommand{\rw}{\rightarrow}
\newcommand{\rrw}{\rightrightarrows}
\newcommand{\Rw}{\Rightarrow}
\newcommand{\lw}{\leftarrow}
\newcommand{\Lw}{\Leftarrow}
\newcommand{\lrw}{\leftrightarrow}
\newcommand{\LRw}{\Leftrightarrow}

\newcommand*\productop{\mathbin{\Pi}}
\newcommand{\dprod}{\productop}

% \newcommand{\dcup}{\sqcup}
% \newcommand{\Dcup}{\bigsqcup}

\newcommand{\dcup}{\coprod}
\newcommand{\Dcup}{\coprod}
\newcommand{\set}[2]{\left\{#1 \ \middle|\ #2\right\}}

\newcommand\mapsfrom{\mathrel{\reflectbox{\ensuremath{\mapsto}}}}

\def\rddots{\mathstrut^{.^{.^{.}}}}

\makeatletter
\newcommand*\bigcdot{\mathpalette\bigcdot@{.5}}
\newcommand*\bigcdot@[2]{\mathbin{\vcenter{\hbox{\scalebox{#2}{$\m@th#1\bullet$}}}}}
\makeatother

%% new delimiter
\DeclarePairedDelimiter{\ceil}{\lceil}{\rceil}
% \emptyset
% \underline

%  ╔══════════════════════════════════════════════════════════╗
%  ║                    Custom Commands                       ║
%  ╚══════════════════════════════════════════════════════════╝

% swap varphi and phi
\let\temp\phi
\let\phi\varphi
\let\varphi\temp



%mod should not have the crazy amount of space to its right!
\renewcommand{\mod}{\,\operatorname{mod}\,}


\newcommand{\nbot}{\mathrel{\rlap{$\bot$}\mkern2mu/}}

%a nice large circle
\newcommand{\tikzcircle}[2][red,fill=black]{\tikz[baseline=-0.5ex]\draw[#1,radius=#2] (0,0) circle ;}%



%  ╔══════════════════════════════════════════════════════════╗
%  ║                     environments                         ║
%  ╚══════════════════════════════════════════════════════════╝


%remember the option: breakable

% create tcolor boxes
\newtcbtheorem[number within=section]{thm}{Theorem}%
{
colback=green!5,
colframe=green!35!black,
fonttitle=\bfseries,
label type=theorem
}{th}

\newtcbtheorem[number within=section, use counter from=thm]{lem}{Lemma}%
{
colback=blue!5,
colframe=black!35!black,
fonttitle=\bfseries,
label type=lemma
}{lm}
\newtcbtheorem[number within=section, use counter from=thm]{prop}{Proposition}%
{
colback=white!5,
colframe=white!35!black,
fonttitle=\bfseries,
label type=proposition
}{pr}
\newtcbtheorem[number within=section, use counter from=thm]{cor}{Corollary}%
{colback=blue!5,
colframe=blue!35!black,
fonttitle=\bfseries,
label type=corollary
}{co}
\newtcbtheorem[number within=section, use counter from=thm]{axiom}{Axiom}%
{colback=black!5,
colframe=yellow!35!black,
fonttitle=\bfseries,
label type=axiom
}{ax}
\newtcbtheorem[number within=section, use counter from=thm]{defn}{Definition}%
{colback=red!5,
colframe=red!35!black,
fonttitle=\bfseries,
label type=definition
}{df}


% for <s-cr> to create \item[$\LRw$] instead of \item
\newenvironment{equivEnumerate}
 {\begin{enumerate}
	}{
\end{enumerate}}



%Custom enviroments
 \newenvironment{note}
 {\begin{tcolorbox}[enhanced, sharp corners, colback=white , borderline={0.2pt}{0pt}{black}]
   \begin{quote}\textbf{Note}:
   }{
   \end{quote}\end{tcolorbox}}

 \newenvironment{intuition}
 {\begin{quote}\textbf{Intuition}:
   }{
   \end{quote}}

\newtcbtheorem[number within=section, use counter from=thm]{titledBox}{}{%
  enhanced,
  sharp corners,
  colback=white,
  colframe=black,
  borderline={0.2pt}{0pt}{black},
  left=2mm,
  right=2mm,
  top=2mm,
  bottom=2mm,
  title style={opacity=1},
  attach title to upper,
  before upper={\textbf{\tcbtitletext}\quad},
  coltitle=black,
  fonttitle=\bfseries,
  separator sign={\quad}
}{box}


%better example and proof environments
\def\exampletext{Example} % If English
\colorlet{colexam}{red!55!black} % Global example color


\newtcbtheorem[number within=section, use counter from=thm]{example}{Example}{% \newtcolorbox[use counter=testexample]{testexamplebox}{%
% Example Frame Start
empty,% Empty previously set parameters
title={\exampletext \thetcbcounter #1},% use \thetcbcounter to access the testexample counter text
% Attaching a box requires an overlay
attach boxed title to top left,
   % Ensures proper line breaking in longer titles
   minipage boxed title,
% (boxed title style requires an overlay)
boxed title style={empty,size=minimal,toprule=0pt,top=4pt,left=3mm,overlay={}},
coltitle=colexam,fonttitle=\bfseries,
before=\par\medskip\noindent,parbox=false,boxsep=0pt,left=3mm,right=0mm,top=2pt,breakable,pad at break=0mm,
   before upper=\csname @totalleftmargin\endcsname0pt, % Use instead of parbox=true. This ensures parskip is inherited by box.
% Handles box when it exists on one page only
overlay unbroken={\draw[colexam,line width=.5pt] ([xshift=-0pt]title.north west) -- ([xshift=-0pt]frame.south west); },
% Handles multipage box: first page
overlay first={\draw[colexam,line width=.5pt] ([xshift=-0pt]title.north west) -- ([xshift=-0pt]frame.south west); },
% Handles multipage box: middle page
overlay middle={\draw[colexam,line width=.5pt] ([xshift=-0pt]frame.north west) -- ([xshift=-0pt]frame.south west); },
% Handles multipage box: last page
overlay last={\draw[colexam,line width=.5pt] ([xshift=-0pt]frame.north west) -- ([xshift=-0pt]frame.south west); }%
}{ex}



\def\prooftext{Proof} % If English
\NewDocumentEnvironment{Proof}{ O{} }
{
    \colorlet{colexam}{black!55!black} % Global example color
    \newtcolorbox[number within=section]{testproofbox}{%
        empty,% Empty previously set parameters
        title={\emph{\prooftext} \thetcbcounter: #1},
        attach boxed title to top left,
        minipage boxed title,
        boxed title style={empty,size=minimal,toprule=0pt,top=4pt,left=3mm,overlay={}},
        coltitle=colexam,
        fonttitle=\bfseries,
        before=\par\medskip\noindent,
        parbox=false,
        boxsep=0pt,
        left=3mm,
        right=0mm,
        top=2pt,
        breakable,
        pad at break=0mm,
        before upper=\csname @totalleftmargin\endcsname0pt,
        % Removed height constraints
        % Added flexible width
        width=\dimexpr\textwidth-3mm\relax,
        % Modified overlays
        overlay unbroken={\draw[colexam,line width=.5pt] ([xshift=-0pt]title.north west) -- ([xshift=-0pt]frame.south west); },
        overlay first={\draw[colexam,line width=.5pt] ([xshift=-0pt]title.north west) -- ([xshift=-0pt]frame.south west); },
        overlay middle={\draw[colexam,line width=.5pt] ([xshift=-0pt]frame.north west) -- ([xshift=-0pt]frame.south west); },
        overlay last={\draw[colexam,line width=.5pt] ([xshift=-0pt]frame.north west) -- ([xshift=-0pt]frame.south west); },%
    }
    \begin{testproofbox}
}
{\end{testproofbox}\endlist}


%  ╔══════════════════════════════════════════════════════════╗
%  ║                    for Dynkin diagrams                   ║
%  ╚══════════════════════════════════════════════════════════╝

\def\row#1/#2!{#1_{\IfStrEq{#2}{}{n}{#2}} & \dynkin{#1}{#2}\\}
\newcommand{\tble}[1]{
   \renewcommand*\do[1]{\row##1!}
   \[
      \begin{array}{ll}\docsvlist{#1}\end{array}
   \]
}

%  ╔══════════════════════════════════════════════════════════╗
%  ║                         links                            ║
%  ╚══════════════════════════════════════════════════════════╝

\definecolor{LinkColour}{HTML}{A64A3B} % favourite
% \definecolor{LinkColour}{HTML}{8C3D32} % 2nd place
\hypersetup{
  colorlinks,
  citecolor=black,
  filecolor=magenta,
  linkcolor=LinkColour,
  urlcolor=blue,
}


%  ╔══════════════════════════════════════════════════════════╗
%  ║                          misc                            ║
%  ╚══════════════════════════════════════════════════════════╝

\stackMath %to be able to stack text in math mode
\makeindex

%import options: all, bibliography, book-formatting, booksSetting, customEnvironments, dynkin, hw-formatting, language-setup, newcommands, packages, preamble, proofs, settings, tcolorbox, test.svg,
\begin{document}
\maketitle


Consider the heat equation on a compact interval:
\begin{equation}\label{eq:linHeatEq}
	\frac{du}{dt} = \frac{\partial^2 u}{\partial x^2}
\end{equation}
where $u$ is a solution. If $u$ represents a rod that is at time $t = 0$ half hot and half cold and is touching in the
middle, then we let time move forward, the temperature will dampen out due to the 2nd term. The shape goes from
a bump function, to a function that looks like $\tan^{-1}(x)$, and as time $\to \infty$ it converges to
a constant function representing the average temperature. These are all functions with distinctly different behaviours.

Now, assume that the heat of the rod at $t = 0$ has the shape of $\sin(x)$ or $\cos(x)$. Then as we let time
propagate these functions dampen uniformly, notably \emph{they are always
the same function}, simply scaled. This is essentially the behaviour of \emph{eigenvectors}, namely as time
moves forward we are scaling by a factor. To have eigenvectors, we need a linear transformation between vector
spaces. Let $S_t$ be the \emph{time evolution operator} that takes in a state $f_{t_0}$ (a function on the
domain that is a $t_0$-slice of a solution $u(x,t_0)$) and outputs $f_{t_0+t}$, in particular $S_t: \SP(X) \to
\SP(X)$ where $\SP(X)$ is the state space over $X$  and $S_t$ is the evolve by $t$ operator:
\[
	S_t(u(-, t_0)) = u(-, t_0+t)
\]
Or, viewed as an (abelian) monoid action\footnote{The heat equation is not well-defined on the full real line
$\R$, it fails the Hadamard’s conditions for well-posedness; see~\cite{NatePDE}}:
\[
	S: [0, \infty) \times \SP(X) \to \SP(X) \qquad t\cdot f_{t_0} = f_{t_0+t}
\]
Note that it isn't ``obvious'' what type of space $\SP(X)$ is, we will harmlessly assume that this is the
Hilbert spaces; see~\cite{NatePDE} for the justification of this assumption. At minimum it is certainly the
case that any (smooth) $L^2$-integrable function is a valid starting state that can then evolve according to
the heat equation. If the reader can accept that there is a reasonable sense in which integrable function can
be differentiated as shall be demonstrated~\cite{NateRealAnalysis}, the reader can let $\SP(X)
= L^2([0,1])$.

Now, certainly $S_t$ is linear as the heat equation is linear and homogeneous\footnote{It is also continuous and
normal, but for that we would need to choose whether we consider $\SP$ as a Banach space or a distribution
space; intuitively the heat equation dissipates energy so $\|S_t f\| \le \|f\|$}. We can thus ask if it has
eigenvectors, that is function $f_r \in \SP(X)$ where
\[
	S_tf_r(x) = \lambda f_{r+t}(x)
\]
If we let $t$ vary, we can define a function $\lambda(t)$ dependent on the time. Given a compact domain $X
=[a,b]$, by direct computations we can check that $\sin(kx)$ is an eigenvector, where this is indeed a state
as it is a $t$-slice of:
\[
	u(x,t) = e^{-k^2 t} \sin(kx)
\]
For example at $t = 0$:
\[
S_t(\sin(kx)) = e^{-k^2 t} \sin(kx)
\]
With the assumption that $\SP$ is a Hilbert space, if these eigenfunctions form an orthonormal basis, we can
decompose a functions into a ``weighted sum'' of these well-behaved functions with respect to time. For the
rest of this section we shall take for granted $\{\sin(kx)\}_{k \in \N}$ form orthonormal basis for $S_t$
(proven in~\cite{NateRealAnalysis}) after appropriately normalized depending on the
domain.

It must be pointed out that the same argument works for $\cos(kx)$:
\[
S_t(\cos(kx)) = e^{-k^2 t} \cos(kx)
\]
and these \emph{also} form an orthonormal basis that \emph{also} diagonalize $S_t$. These two sets certainly
cannot be combined since the result is no longer orthogonal:
\[
	 \int_0^1 \sin(\pi x) \cos(2\pi x)  dx = \frac{-2}{3\pi} \ne 0
\]
Thus, the question shifts into one about finding a \emph{natural} eigenbasis for $S_t$. A-priori, it is not
necessarily the case that these are the only eigenbasis. What we need is for the eigenbasis to ``respect''
some inherent geometry of the space. Phrased differently, the eigenbasis needs to be \emph{invariant} with
respect to some operator(s), where the operators represent some symmetry of the space. To get a better grasp
of invariance of operators, we may start looking at other PDEs and considering diagonalizing the operator
$S_t$ with respect to the solutions of those PDEs.


For example, if we solve the Heat equation on a sphere:
\[
	\frac{\partial u}{\partial t} = \Delta_{\text{sphere}} u
\]
the operator can still be diagonalized (see~\cite{NatePDE}). What we shall want is for an eigenbasis that
respects the spherical geometry,
namely for it to be invariant to an $SO(3)$ action. Such eigenfunctions do exist, they are the \emph{Spherical Harmonics}
$Y_{l}^{m}(\theta, \phi)$, and they are invariant to rotation\footnote{important technicality: $SO(3)$ is not
abelian, hence it cannot be that \emph{all} rotation operators simultaneously commute. The best solution
is to pick a ``maximal set'' of commuting symmetries: usually the Total Angular Momentum ($L^2$,
representing the amount of rotation invariance) and Azimuthal Symmetry ($L_z$, rotation around the
z-axis).}
\[
	\Delta (f \circ R) = (\Delta f) \circ R
\]
To give another example that is famous in physics, consider Quantum Harmonic Oscillator
\[
	i \frac{\partial \psi}{\partial t} = -\frac{\partial^2 \psi}{\partial x^2} + x^2 \psi
\]
This PDE is diagonalizable (see~\cite{NatePDE}), but it is \emph{not} diagonalizable by $e^{2\pi i k}$.
Instead it was shown to be diagonalized (uniquely) by the \emph{Hermite Functions} (a Gaussian multiplied by a polynomial)
\[
	\phi_n(x) = H_n(x) e^{-x^2/2}
\]
By trying to diagonalize with $e^{2\pi k}$, the $x^2$ multiplication term turns into a Laplacian in $k$-space
($-\partial_k^2$), so the result will effectively just swap the PDE for an identical PDE in frequency space
without solving it. In physics, this fact shall be the key to showing that it is in fact \emph{invariant} to
the Fourier transform, creating fascinating link between abstract harmonic analysis and the concrete geometry
of quantum mechanics, a connection to be explored in another book.

% Another importance invariance that should be mentioned for its use in number theory is \emph{dilation
% invariance} which results in the \emph{Mellin transform}, but it has no bearing on this section so



So what is the invariance we want to impose for the operator $S_t$ with respect to the heat equation? In the
example of the sphere, we had that the group $SO(3)$. A compact interval has no invariance simply because it
is not a group. However, by gluing the end points of the interval we get a 1-\emph{torus} (a
circle)\footnote{More generally, any rectangular domain is homeomorphic to $[0,1]^n$ which then can be glued
to become $\T^n = (S^1)^n$, and so this same translation structure appears in higher dimensions}. We now have
the rotation group operation on $\T^1$, or equivalently the translation operation on $\R/\Z$\footnote{Really
$\C/\Z$ as we shall justify shortly}. This does shift the geometry and hence the inherent dynamics of our
PDEs; the problem of solving a PDE on a compact interval will be a subset of solving PDEs on a torus as shall
be addressed shortly. Let us first simply consider $S_t$ on $L^2(\T^1)$ and see how it interacts with the
translation operator. Since we are working on the torus, we now have that $\sin(2\pi kx)$ and $\cos(2\pi lx)$
\emph{are} orthogonal (where the $2\pi$-scaling was to ensure a full period of each lies within the torus),
and $\{\sin(kx), \cos(lx)\}_{k,l \in \N}$ is an orthonormal basis (after appropriate scaling and stretching).
Now, note that $\sin$ and $\cos$ basis are \emph{not} translation invariant, for example:
\[
	\tau_a(\sin(kx)) = \cos(ka)\sin(kx) + \sin(ka)\cos(kx)
\]
and so they are not eigenvectors of $\tau_a$. However, they still \emph{block-diagonalize} the translation
operator $\tau_a$ with blocks of size $2$:
\[
	\begin{pmatrix} \cos(ka) & \sin(ka)\\
	 -\sin(ka) & \cos(ka) \end{pmatrix}
\]
If we look at these blocks, we shall see that they are all rotation matrices by an angle $ka$; in other words the
imposition is that we did not have enough eigenvalues. By extending the codomain to $\C$, we can readily
verify that:
\[
	\tau_a(e^{iz}) = e^{ia}\cdot e^{iz}
\]
and so $\{e^{ikz}\}_{k \in \Z}$ is our desired eigenbasis for $\tau_a$. Now, do these translations of these
functions commute with $S_t$, namely:
\[
	S_t(\tau_a(f)) = \tau_a(S_t(f)) \qquad \forall a \in \T^1
\]
Indeed:
\begin{align*}
	\tau_a(S_t(e^{ikx})) &= \tau_a(e^{-k^2 t} e^{ikx})\\
	&=  e^{-k^2 t} e^{ik(x+a)}\\
	&= e^{-k^2 t} e^{ika} e^{ikx}\\
	&= e^{ika} S_t(e^{ikx})\\
	&= S_t(e^{ika} e^{ikx}) & \text{$S_t$ is linear}\\
	&= S_t(\tau_a(e^{ikx}))
\end{align*}

What this means physically is that the heat equation dissipates the same way no matter where you are on a
ringed-shape rod. If the rod was half wood and half metal, then this translation invariance shall fail.

Now, notice that as $\tau_a \circ S_t = S_t \circ \tau_a$ for all $a \in \T^1$ (or in commutator bracket
notation $[\gamma_a, S_t] = 0$ for all $a,t$), the eigenvectors of $\tau_a$ are the same eigenvectors as
$S_t$. The eigenvectors of $S_t$ are in general hard to compute, but the eigenvectors of $\tau_a$ are
\emph{precisely} the functions $e^{ikx}$.


\subsection*{Boundary Problem}

The immediate point that is still pending and must clarified when switching to $\T^1$ is that this shifts the
dynamics of any PDE we are solving, as now points at the boundary are affected by their neighbourhood.
However, for PDE's that will respect the group structure of $\T^1$, the effects at the boundary can be
``cancelled out'' in two separate ways that shall recover the original two basis $\{\sin(kx)\}_{k \in \Z}$ and
$\{\cos(kx)\}_{k \in \Z}$ we saw before (again, appropriately normalized). The key insight is that the
interval $[0,1]$ can be viewed as a \emph{sub-domain} of the torus $\T^1$ (identified with $[-1,1]$) subject
to a symmetry constraint. While the torus possesses full translation invariance, the interval possesses
\emph{reflection invariance}. By imposing a symmetry on the torus, we create ``virtual boundaries'' at the
fixed points of the reflection ($x=0$ and $x=1$). This gives rise to two distinct orthonormal bases for
$L^2([0,1])$, each corresponding to a different physical boundary condition:

\begin{itemize}
    \item \emph{Odd Symmetry (Dirichlet Conditions):}
    Consider the subspace of functions on $\T^1$ that are \emph{odd}, satisfying $u(-x) = -u(x)$. At the
    fixed point $x=0$, this symmetry forces $u(0) = -u(0)$, implying $u(0) = 0$. Physically, the heat
    flowing from the right is perfectly cancelled by the ``negative heat'' flowing from the left, creating a
    virtual cold wall. The eigenfunctions respecting this symmetry are precisely the sine functions:
    \[ \sin(n\pi x) = \frac{e^{in\pi x} - e^{-in\pi x}}{2i} \]
    Restricting these functions to $[0,1]$ yields the \emph{Sine Basis}, which diagonalizes the Laplacian
    subject to $u(0)=u(1)=0$.

    \item \emph{Even Symmetry (Neumann Conditions):}
    Consider the subspace of functions on $\T^1$ that are \emph{even}, satisfying $u(-x) = u(x)$. At the
    fixed point, the derivative must be odd, forcing $u'(0) = 0$. Physically, this represents an insulated
    boundary where heat reflects back without loss. The eigenfunctions respecting this symmetry are the cosine
    functions:
    \[ \cos(n\pi x) = \frac{e^{in\pi x} + e^{-in\pi x}}{2} \]
    Restricting these to $[0,1]$ yields the \emph{Cosine Basis}, which diagonalizes the Laplacian subject
    to $u'(0)=u'(1)=0$.
\end{itemize}

This leads to a remarkable result in Harmonic Analysis known as \emph{half-range expansions}: while sines and
cosines \emph{together} form a basis for the full torus, \emph{either} set alone forms a complete orthonormal
basis for the interval $[0,1]$.





\subsection*{Decoupling PDE}

To understand the consequences better, take a solution $u: \T^1\times [0, \infty) \to \C$. Then, think of
each (space) neighborhood in the usual heat equation as influencing its infinitesimal neighbourhood, namely
the $\frac{\partial^2 u}{\partial x^2}$ term determines how each neighbouring point affects each other as
time propagates. We can think of this as being uncountably many dependent equations.

Therefore, by applying the Fourier transform, we are changing the basis of our state space
from the position basis $\{ \delta_x \}_{x \in \T^1}$ to the eigenbasis of the spatial
operator $\{ e^{ikx} \}_{k \in \Z}$.

By taking the Fourier transform, i.e. transforming to the
translation-invariant orthonormal basis, we change a solution to:
\[
	u(x,t): \T^1 \times [0, \infty) \to \C\qquad \xmapsto{\CF_x}\qquad \hat{u}(k, t): \Z\times [0,\infty) \to \C
\]
where the $\Z$ is the index for the eigenbasis $\{f_k\}_{k \in \Z}$. Then if $\hat{u}_k$ is the function whose
input is the index of the eigenvector and output is its coefficient, we see that as $\hat{u}_k$ propagates
through time, $\hat{u}_k(t)$ is only dependent on the original function up to scaling. In fact
we shall show that the heat equation PDE decomposes into the equations:
\[
	\frac{d \hat{u}_k}{dt} = -k^2 \hat{u}_k \qquad \forall k \in \Z
\]
where now we have \emph{countably many decoupled} equations.


\subsection*{Non-Compact Domain}

Notice that $\sin, \cos, e^{z}$ are not $L^2$-integrable on $\R^n$ (resp. $\C^n$), hence these cannot form an
orthonormal basis on $L^2(\R^n)$. Thus, the initial foray into Fourier will focus on functions with domain
$\T^n$. In~\cite{NateRealAnalysis} we shall show how to generalize the tools to the non-compact
case, more specifically, the locally compact case.


\subsection*{Digression to Lie Theory}

We have established that the Fourier basis functions $e^{ikx}$ are the eigenfunctions of the ``global''
spatial translation operator $\tau_a$. However, the reader likely recognizes these functions primarily as
the solutions to the differential equation $f' = \lambda f$. This connection between the \emph{global}
translation $\tau_a$ and the \emph{local} derivative $D = \frac{d}{dx}$ is a central theme of Lie Theory.

In~\cite{NateRealAnalysis}, it shall be shown that it makes sense to infinitely differentiate
integrable functions. For now, we rely on the intuition that any continuous function can be approximated by
a smooth function (and locally by a polynomial).

Consider the Taylor expansion of a function $f$ around a point $x$:
\[
	f(z) = f(x) + f'(x)(z-x) + \frac{f''(x)}{2!}(z-x)^2 + \dots
\]
If we wish to evaluate $f$ at a shifted point $x+y$, the Taylor series gives:
\[
	f(x+y) = f(x) + y f'(x) + \frac{y^2}{2!} f''(x) + \dots
\]
Notice that we can re-write this using the differentiation operator $D = \frac{d}{dx}$:
\[
	f(x+y) = \left( I + yD + \frac{y^2 D^2}{2!} + \dots \right) f(x)
\]
Recognizing the series inside the parenthesis as the definition of the matrix exponential, we arrive at the
compact operator identity:
\[
	\tau_y = e^{y D}
\]
This equation states that \emph{differentiation generates translation}.

Now, consider the spectral consequences. If an operator $A$ has an eigenvector $v$ with eigenvalue $\lambda$
($Av = \lambda v$), then the operator exponential $e^A$ acts on $v$ as:
\[
e^A v = \left( I + A + \frac{A^2}{2!} + \dots \right) v = \left( 1 + \lambda + \frac{\lambda^2}{2!} + \dots \right) v = e^\lambda v
\]
Applying this to our spatial operators: since $f(x) = e^{ikx}$ is an eigenfunction of the derivative $D$
with eigenvalue $ik$ (i.e., $D e^{ikx} = ik e^{ikx}$), it must automatically be an eigenfunction of the
translation operator $\tau_y = e^{yD}$ with eigenvalue $e^{iyk}$.

This confirms our earlier finding: the Fourier basis $e^{ikx}$ simultaneously diagonalizes the derivative
(and by extension the Laplacian $\Delta = D^2$) and the translation group.

\begin{thm}{Spectral Mapping Theorem}{spectralMappingThm}
	Any eigenvector with eigenvalue $\lambda$ of an operator $A$ is an eigenvector of the operator $e^A$ with
	eigenvalue $e^\lambda$.
\end{thm}

\begin{Proof}
	Follows from the power series definition of the matrix exponential shown above.
\end{Proof}

To see this explicitly, observe the action of these operators on the basis function $f_k(x) = e^{ikx}$:
\begin{itemize}
	\item \emph{Derivative:} $D f_k = (ik) f_k$ \hfill (Eigenvalue $ik$)
	\item \emph{Laplacian:} $\Delta f_k = D^2 f_k = (ik)^2 f_k = -k^2 f_k$ \hfill (Eigenvalue $-k^2$)
	\item \emph{Translation:} $\tau_a f_k = e^{aD} f_k = e^{a(ik)} f_k$ \hfill (Eigenvalue $e^{ika}$)
\end{itemize}

In the language of representation theory, the differentiation operator is the \emph{Lie Algebra} generator
$\mathfrak{g}$, and the translation operator is the \emph{Lie Group} element $e^{\mathfrak{g}}$. The Fourier
basis provides the unique coordinate system where the group action and its generator are both simple scalings.

\subsection*{Generalizing using Representation Theory}

In the previous section, our motivation was to solve a specific time-dependent PDE (the heat equation). We
discovered that the key to solving it lay not in the time evolution $S_t$ itself, but in the underlying
symmetry of the spatial domain (translation invariance). The basis functions $e^{ikx}$ worked because they
diagonalized the translation operator $\tau_a$.

We now abstract this principle. We no longer need a specific ``time'' equation; instead, we look for the
basis that diagonalizes the action of the symmetry group itself. By finding the irreducible representations
of the symmetry group, we construct a ``universal'' eigenbasis that will automatically diagonalize \emph{any}
linear operator respecting that symmetry (be it the Laplacian, the Heat Operator, or others).

To make this concrete, let us recall the representation theory of finite groups.
If $G$ is a finite group, we can represent $G$ on a vector space $V$ (over a field $k$) by
taking a homomorphism $\rho: G \to \GL(V)$.
\[
    G \actson V \qquad \rho(g)v = g \cdot v
\]
As we are not interested in any arithmetic limitations, we take $k = \C$. Then by Maschke's theorem
(see~\cite[chapter~24]{NateAlgebra}) the induced $\C[G]$-module structure on $V$ is
semi-simple, meaning:
\[
	 V \cong_G \bigoplus_i V_i
\]
where $V_i$ are irreducible representations. Given the basis for the above isomorphism, any $g \in G$
will give a matrix representation $\rho(g)$ which respects this block-diagonal decomposition:
\[
\rho(g) = \begin{bmatrix}
\mathbf{M}_1(g) & 0 & 0 & 0 \\
0 & \mathbf{M}_2(g) & 0 & 0 \\
0 & 0 & \mathbf{M}_3(g) & 0 \\
0 & 0 & 0 & \dots
\end{bmatrix}
\]
where $\mathbf{M}_i(g)$ represents the matrix block acting on the subspace $V_i$. Crucially, if $G$ is
\emph{abelian}, Schur's Lemma implies that \emph{all} irreducible representations are $1$-dimensional. Hence,
the matrices become scalars:
\[
\rho(g) = \begin{bmatrix}
\lambda_1(g) & 0 & 0 & 0 \\
0 & \lambda_2(g) & 0 & 0 \\
0 & 0 & \lambda_3(g) & 0 \\
0 & 0 & 0 & \dots
\end{bmatrix}
\]
As each $\rho(g)$ is representable by a diagonal matrix using the same basis, the operators $\{\rho(g)\}_{g
\in G}$ are a commuting family of linear operators (i.e. they are \emph{simultaneously diagonalizable}).

Now, we generalize this to function spaces. While for a finite group the vector space of functions $\C[G]$ is
finite-dimensional, equipping it with an inner product structure is essential to define unitarity and
orthogonality. A finite group $G$ carries a natural unique measure (the counting measure), which we can
normalize to be a probability measure (the Haar measure $\mu(g) = 1/|G|$). This allows us to define the
Hilbert space $L^2(G)$ with inner product:
\[
    \langle f, h \rangle = \frac{1}{|G|} \sum_{x \in G} f(x)\overline{h(x)}
\]
We define the \emph{Right Regular Representation}\footnote{or, Left Regular
Representations $(L_g\cdot f)(x) = f(g^{-1}x)$} of $G$ on this space:
\[
	G \actson L^2(G) \qquad (R_g \cdot f)(x) = f(xg)
\]
We verify that this action is unitary (and hence can be diagonalized with an orthonormal basis). Since
right-multiplication by $g$ is a bijection on the group $G$ (it permutes the elements), the sum over $x$ is
identical to the sum over $y = xg$:
\begin{align*}
    \langle R_g f, R_g h \rangle &= \frac{1}{|G|} \sum_{x \in G} f(xg)\overline{h(xg)} \\
    &= \frac{1}{|G|} \sum_{y \in G} f(y)\overline{h(y)} \\
    &= \langle f, h \rangle
\end{align*}
Thus, the operators $R_g$ are unitary. Since $G$ is abelian, the operators $\{R_g\}_{g \in G}$ form a
commuting family of unitary operators. By the spectral theorem for unitary operators, there exists an
orthonormal basis $\mathcal{B}$ of $L^2(G)$ that simultaneously diagonalizes the action.
If $\phi \in \mathcal{B}$ is such a basis vector, then:
\[
	R_g(\phi) = \chi_{\phi}(g)\phi
\]
where $\chi_{\phi}(g) \in \C$ is the eigenvalue. Since the representation is a homomorphism ($R_{gh} = R_g
R_h$), the eigenvalues must satisfy the multiplicative property $\chi_\phi(gh) = \chi_\phi(g)\chi_\phi(h)$.
Moreover, since $R_g$ is unitary, $|\chi_\phi(g)|=1$ and hence the codomain is $S^1$. These homomorphisms
$\chi: G \to S^1$ are precisely the \emph{characters} of the group. Thus,

\begin{quote}
	the group-theoretic eigenbasis for $L^2(G)$ is simply the set of characters!
\end{quote}

We shall explore this more in~\cite{NateRealAnalysis}

\begin{titledBox}{Advanced: Operators Commuting with Translations}{advancedOpCommutingTrans}
    We focused on finding the basis that diagonalizes translations ($R_g$). But what are the operators that
    commute with all translations?
    \[ T \circ R_g = R_g \circ T \qquad \forall g \in G \]
    A fundamental result in Harmonic Analysis states that any such linear operator $T$ is a \emph{Convolution Operator}.
    That is, there exists a function $k \in L^1(G)$ (the kernel) such that $Tf = f * k$.
    This explains why the Heat Operator, the Laplacian, and other physical operators can all be written as convolutions:
    they are simply the most general linear maps that respect the symmetry of the space.
\end{titledBox}

\subsubsection*{The Infinite Compact Case}

In the case where the domain is infinite (like the torus $\T^n$), we must be more careful. First, let us
stick to compact domains so that we generalize the finite case directly. Functions defined on a rectangular
domain with periodic boundary conditions are equivalent to functions on the torus $\T^n = (S^1)^n$.
Our state space is the Hilbert space $L^2(\T^n, \C)$.

We want to act by translation by $\R^n$, but since the domain is periodic, this descends to a faithful action
of the compact group $\R^n/\Z^n \cong \T^n$. It is easy to check that translation is unitary ($\langle
\tau_y f, \tau_y g \rangle = \langle f,g \rangle$) using the translation invariance of the Lebesgue measure.

We now invoke the \emph{Peter-Weyl Theorem}, which is the generalization of Maschke's theorem to compact
topological groups (see~\cite{NateGeoRepTheory}). It states that the unitary representation of a compact
group on $L^2(G)$ decomposes into a Hilbert direct sum of finite-dimensional irreducible representations:
\[
	L^2(G) \cong_{G} \widehat{\bigoplus_i} V_i
\]
Since $\T^n$ is abelian, all irreducible representations $V_i$ are \emph{one-dimensional}. Being one-dimensional
and unitary, they are given by characters $\chi: \T^n \to S^1$.

Thus, the Peter-Weyl theorem guarantees that the characters of $\T^n$ form a complete orthonormal basis for
$L^2(\T^n)$. We have recovered our Fourier basis purely from the algebraic structure of the domain, without
reference to any specific differential equation!

\subsection*{* The Duality of Points and Spectrum}


The statement that the Fourier transform changes the basis from ``positions'' $\{\delta_x\}$ to ``waves''
$\{e^{ikx}\}$ hints at a broader duality: it changes the very nature of the ``points'' that constitute our
space.

Let us define a function $f$ by its values at geometric points. In the language of distributions, we can write
any function as a continuous sum of Dirac deltas:
\[
    f(x) = \int_{\T^1} f(y) \delta_y(x) \, dy
\]
Here, the building blocks are points $y \in \T^1$. A ``point'' is simply a linear functional that evaluates a
function at a specific location: $\text{ev}_y(f) = f(y)$.

Now consider the Fourier Transform. It maps a function $f$ on the group $G$ to a function $\hat{f}$ on the
dual group $\widehat{G}$ (the set of characters). What are the ``points'' of this new space? They are the
irreducible representations (the frequencies).
\[
    \hat{f}(\chi) = \int_G f(x) \overline{\chi(x)} \, dx
\]
When we transform a geometric point $\delta_y$, it becomes a function on the dual group:
\[
    \widehat{\delta_{y}}(\chi) = \overline{\chi(y)}
\]
This reveals that the Fourier Transform is an isomorphism between two different kinds of spaces:
\begin{itemize}
    \item The \emph{Physical Space} $G$, whose points are geometric locations $x$.
    \item The \emph{Spectral Space} $\widehat{G}$, whose points are irreducible representations $\chi$.
\end{itemize}

This is a specific instance of a grander theme in mathematics known as \emph{Gelfand Duality}. In this
framework, any commutative $C^*$-algebra (like the algebra of continuous functions) corresponds to a
topological space (its spectrum). The Fourier Transform is the map that identifies the algebra of functions
under convolution ($L^1(G)$) with the algebra of functions under pointwise multiplication ($C_0(\widehat{G})$).
Effectively, we trade the geometry of the group for the geometry of its representation theory.

\newpage
\printbibliography

\end{document}
